% !TEX program = xelatex
% Options for packages loaded elsewhere
\PassOptionsToPackage{unicode}{hyperref}
\PassOptionsToPackage{hyphens}{url}
\PassOptionsToPackage{dvipsnames,svgnames,x11names}{xcolor}
\documentclass[
  8pt,
  twocolumn]{extarticle}
\usepackage{xcolor}
\usepackage[landscape,margin=0.38in]{geometry}
\usepackage{amsmath,amssymb}
\setcounter{secnumdepth}{-\maxdimen} % remove section numbering
\usepackage{iftex}
\ifPDFTeX
  \usepackage[T1]{fontenc}
  \usepackage[utf8]{inputenc}
  \usepackage{textcomp} % provide euro and other symbols
\else % if luatex or xetex
  \usepackage{unicode-math} % this also loads fontspec
  \defaultfontfeatures{Scale=MatchLowercase}
  \defaultfontfeatures[\rmfamily]{Ligatures=TeX,Scale=1}
\fi
\usepackage{lmodern}
\ifPDFTeX\else
  % xetex/luatex font selection
\fi
% Use upquote if available, for straight quotes in verbatim environments
\IfFileExists{upquote.sty}{\usepackage{upquote}}{}
\IfFileExists{microtype.sty}{% use microtype if available
  \usepackage[]{microtype}
  \UseMicrotypeSet[protrusion]{basicmath} % disable protrusion for tt fonts
}{}
\makeatletter
\@ifundefined{KOMAClassName}{% if non-KOMA class
  \IfFileExists{parskip.sty}{%
    \usepackage{parskip}
  }{% else
    \setlength{\parindent}{0pt}
    \setlength{\parskip}{6pt plus 2pt minus 1pt}}
}{% if KOMA class
  \KOMAoptions{parskip=half}}
\makeatother
\setlength{\emergencystretch}{3em} % prevent overfull lines
\providecommand{\tightlist}{%
  \setlength{\itemsep}{0pt}\setlength{\parskip}{0pt}}
\usepackage{microtype}
\usepackage{mathtools}
\usepackage{amssymb}
\setlength{\columnsep}{0.32in}
\setlength{\parindent}{0pt}
\setlength{\parskip}{2pt}
\setlength{\abovedisplayskip}{4pt}
\setlength{\belowdisplayskip}{4pt}
\setlength{\abovedisplayshortskip}{3pt}
\setlength{\belowdisplayshortskip}{3pt}
\allowdisplaybreaks
\sloppy
\emergencystretch=3em
\usepackage{bookmark}
\IfFileExists{xurl.sty}{\usepackage{xurl}}{} % add URL line breaks if available
\urlstyle{same}
\hypersetup{
  colorlinks=true,
  linkcolor={black},
  filecolor={Maroon},
  citecolor={Blue},
  urlcolor={black},
  pdfcreator={LaTeX via pandoc}}

\author{}
\date{}

\begin{document}

\section{Formula Sheet Part 2: Asset Pricing, Diagnostics, and
EMH}\label{formula-sheet-part-2-asset-pricing-diagnostics-and-emh}

This is a two-column topic split from
\texttt{All\_Formulas\_From\_Chatnotes.md}.

\subsection{Conventions}\label{conventions}

\[
t,T,h,k,n,p,q,N
\]

Variables: \(t\) is a time index; \(T\) is sample size or the final
observed time; \(h\) is a forecast horizon; \(k\) is a lag or number of
model parameters depending on context; \(n\) is a return horizon;
\(p,q\) are lag orders; \(N\) is the number of assets.

\[
F_t
\]

Variables: \(F_t\) is the information set available at time \(t\).

Returns may be written as decimals or percentages. Check the scale
before applying VaR or standard deviation formulas.

\subsection{Asset Pricing Models}\label{asset-pricing-models}

\subsubsection{CAPM pricing equation}\label{capm-pricing-equation}

\[
\mathbb E(R_i)-r_f=\beta_i[\mathbb E(R_m)-r_f]
\]

Variables: \(R_i\) is asset return; \(R_m\) is market return;
\(\beta_i\) is market beta.

\subsubsection{CAPM regression}\label{capm-regression}

\[
r_{i,t}-r_{f,t}=\alpha_i+\beta_i(r_{m,t}-r_{f,t})+e_{i,t}
\]

Variables: \(r_{i,t}-r_{f,t}\) is asset excess return;
\(r_{m,t}-r_{f,t}\) is market excess return; \(\alpha_i\) is abnormal
return; \(e_{i,t}\) is the residual.

\subsubsection{CAPM alpha restriction}\label{capm-alpha-restriction}

\[
\alpha_i=0
\]

Variables: \(\alpha_i\) is abnormal return relative to CAPM.

\subsubsection{CAPM total risk
decomposition}\label{capm-total-risk-decomposition}

\[
\operatorname{Var}(r_{i,t}-r_f)=\beta_i^2\operatorname{Var}(r_{m,t}-r_f)+\operatorname{Var}(e_{i,t})
\]

Variables: total excess-return variance equals systematic variance plus
idiosyncratic variance.

\subsubsection{Systematic and idiosyncratic
fractions}\label{systematic-and-idiosyncratic-fractions}

\[
\begin{aligned}
\text{systematic fraction}&=R^2\\
\text{idiosyncratic fraction}&=1-R^2
\end{aligned}
\]

Variables: \(R^2\) is the regression explained-variation statistic.

\subsubsection{Fama-French three-factor
model}\label{fama-french-three-factor-model}

\[
r_{i,t}-r_{f,t}=\alpha_i+\beta_{i1}(r_{m,t}-r_{f,t})+\beta_{i2}SMB_t+\beta_{i3}HML_t+e_{i,t}
\]

Variables: \(\beta_{i1}\) is market loading; \(\beta_{i2}\) is SMB
loading; \(\beta_{i3}\) is HML loading.

\subsubsection{Carhart four-factor
model}\label{carhart-four-factor-model}

\[
r_{i,t}-r_{f,t}=\alpha_i+\beta_{i1}(r_{m,t}-r_{f,t})+\beta_{i2}SMB_t+\beta_{i3}HML_t+\beta_{i4}MOM_t+e_{i,t}
\]

Variables: \(MOM_t\) is the momentum factor; \(\beta_{i4}\) is momentum
loading.

\subsubsection{CAPM nested in the four-factor
model}\label{capm-nested-in-the-four-factor-model}

\[
\beta_{i2}=\beta_{i3}=\beta_{i4}=0
\]

Variables: setting the extra factor loadings to zero reduces the
four-factor model to CAPM.

\subsubsection{Single-coefficient
t-test}\label{single-coefficient-t-test}

\[
\begin{aligned}
H_0&:\beta=\beta_0\\
H_1&:\beta\ne\beta_0\\
t&=\frac{\hat\beta-\beta_0}{\operatorname{se}(\hat\beta)}
\end{aligned}
\]

Variables: \(\hat\beta\) is the estimated coefficient; \(\beta_0\) is
the hypothesised value; \(\operatorname{se}(\hat\beta)\) is the standard
error.

\subsubsection{Test whether a factor
matters}\label{test-whether-a-factor-matters}

\[
\begin{aligned}
H_0&:\beta_{SMB}=0\\
H_1&:\beta_{SMB}\ne0\\
t&=\frac{\hat\beta_{SMB}}{\operatorname{se}(\hat\beta_{SMB})}
\end{aligned}
\]

Variables: \(\beta_{SMB}\) is the SMB coefficient.

\subsubsection{Joint test of CAPM against a multi-factor
model}\label{joint-test-of-capm-against-a-multi-factor-model}

\[
\begin{aligned}
H_0&:\beta_{SMB}=\beta_{HML}=\beta_{MOM}=0\\
H_1&:\text{at least one extra }\beta\ne0\\
J&=\frac{RSS_0-RSS_1}{RSS_1/(T-K-1)}
\end{aligned}
\]

Variables: \(RSS_0\) is restricted-model residual sum of squares;
\(RSS_1\) is unrestricted-model residual sum of squares; \(K\) is the
number of regressors in the unrestricted model.

\subsubsection{Residual standard error and
RSS}\label{residual-standard-error-and-rss}

\[
RSS=(\text{residual standard error})^2df
\]

Variables: \(df\) is residual degrees of freedom.

\[
\text{residual standard error}=\sqrt{\frac{RSS}{df}}
\]

Variables: residual standard error estimates residual standard
deviation.

\[
df=T-K-1
\]

Variables: \(T\) is sample size; \(K\) is the number of regressors
excluding the intercept.

\subsection{Regression Diagnostics and EMH
Tests}\label{regression-diagnostics-and-emh-tests}

\subsubsection{Regression residual
properties}\label{regression-residual-properties}

\[
\begin{aligned}
\mathbb E(e_{i,t})&=0\\
\operatorname{Var}(e_{i,t})&\text{ is constant}\\
\mathbb E(e_{i,t}e_{i,t-k})&=0,\quad k=1,2,\ldots
\end{aligned}
\]

Variables: \(e_{i,t}\) is a regression disturbance or residual.

\subsubsection{Original model and
residual}\label{original-model-and-residual}

\[
y_t=\alpha+\beta x_t+e_t
\]

Variables: \(y_t\) is the dependent variable; \(x_t\) is the regressor;
\(\alpha,\beta\) are coefficients.

\[
\hat e_t=y_t-\hat\alpha-\hat\beta x_t
\]

Variables: \(\hat e_t\) is the fitted residual.

\subsubsection{Heteroskedasticity
hypotheses}\label{heteroskedasticity-hypotheses}

\[
\begin{aligned}
H_0&:\operatorname{Var}(e_{i,t})\text{ is constant}\\
H_1&:\operatorname{Var}(e_{i,t})\text{ is not constant}
\end{aligned}
\]

Variables: the null is homoskedasticity; the alternative is
heteroskedasticity.

\subsubsection{Heteroskedasticity auxiliary
regression}\label{heteroskedasticity-auxiliary-regression}

\[
\hat e_t^2=\gamma_0+\gamma_1x_t+\gamma_2x_t^2+v_t
\]

Variables: \(\gamma_j\) are auxiliary-regression coefficients; \(v_t\)
is the auxiliary error.

\subsubsection{LM statistic for
heteroskedasticity}\label{lm-statistic-for-heteroskedasticity}

\[
W=TR^2
\]

Variables: \(R^2\) comes from the auxiliary regression; compare \(W\)
with a chi-square critical value.

\subsubsection{Residual autocorrelation
hypotheses}\label{residual-autocorrelation-hypotheses}

\[
\begin{aligned}
H_0&:\text{no autocorrelation}\\
H_1&:\text{at least one lagged residual matters}
\end{aligned}
\]

Variables: the alternative means residual serial correlation remains.

\subsubsection{Residual autocorrelation auxiliary
regression}\label{residual-autocorrelation-auxiliary-regression}

\[
\hat e_t=\gamma_0+\gamma_1\hat e_{t-1}+\gamma_2\hat e_{t-2}+\cdots+v_t
\]

Variables: \(\hat e_{t-j}\) are lagged residuals.

\subsubsection{LM statistic for
autocorrelation}\label{lm-statistic-for-autocorrelation}

\[
AR(p)=TR^2
\]

Variables: \(p\) is the number of residual lags; under the null compare
to \(\chi_p^2\).

\subsubsection{HAC definition}\label{hac-definition}

\[
HAC=\text{heteroskedasticity and autocorrelation consistent}
\]

Variables: HAC/Newey-West changes standard errors, not fitted
coefficients.

\subsubsection{ETF tracking regression}\label{etf-tracking-regression}

\[
r_{\text{ETF},t}=\alpha+\beta r_{\text{index},t}+e_t
\]

Variables: \(\beta\) measures the ETF's exposure to the index.

\subsubsection{Leveraged ETF target
test}\label{leveraged-etf-target-test}

\[
\begin{aligned}
H_0&:\beta=-2\\
H_1&:\beta\ne-2\\
t&=\frac{\hat\beta-(-2)}{\operatorname{se}(\hat\beta)}
\end{aligned}
\]

Variables: \(\hat\beta\) is the estimated tracking slope.

\subsubsection{EMH return model}\label{emh-return-model}

\[
r_t=\mu+e_t
\]

Variables: \(\mu\) is constant mean return; \(e_t\) is the unpredictable
return component.

\subsubsection{White noise notation}\label{white-noise-notation}

\[
e_t\sim WN(0,\sigma^2)
\]

Variables: \(WN(0,\sigma^2)\) means white noise with zero mean and
variance \(\sigma^2\).

\subsubsection{White noise conditions}\label{white-noise-conditions}

\[
\begin{aligned}
\mathbb E(e_t)&=0\\
\operatorname{Var}(e_t)&=\sigma^2\\
\operatorname{Corr}(e_t,e_{t-k})&=0,\quad k\ge1
\end{aligned}
\]

Variables: no autocorrelation means past shocks do not linearly predict
current shocks.

\subsubsection{Autocorrelation}\label{autocorrelation}

\[
\rho(k)=\operatorname{Corr}(r_t,r_{t-k})
\]

Variables: \(\rho(k)\) is autocorrelation at lag \(k\).

\subsubsection{Correlogram confidence
band}\label{correlogram-confidence-band}

\[
\pm\frac{2}{\sqrt T}
\]

Variables: \(T\) is sample size; the band is an approximate 5 percent
individual significance guide.

\subsubsection{Q-test hypotheses}\label{q-test-hypotheses}

\[
\begin{aligned}
H_0&:\rho_1=\rho_2=\cdots=\rho_k=0\\
H_1&:\text{at least one autocorrelation is non-zero}
\end{aligned}
\]

Variables: \(\rho_j\) is autocorrelation at lag \(j\).

\subsubsection{Portmanteau/Ljung-Box statistic from the practice
exam}\label{portmanteauljung-box-statistic-from-the-practice-exam}

\[
Q=T(T+2)\sum_{k=1}^{m}\frac{\hat\rho_k^2}{T-k}
\]

Variables: \(m\) is the largest lag tested; \(\hat\rho_k\) is the sample
autocorrelation at lag \(k\); compare \(Q\) with \(\chi_m^2\).

\subsubsection{Single autocorrelation large-sample
test}\label{single-autocorrelation-large-sample-test}

\[
\sqrt T\,\hat\rho(k)\approx\mathcal N(0,1)
\]

Variables: \(\hat\rho(k)\) is the sample autocorrelation at lag \(k\).

\[
\begin{aligned}
H_0&:\rho(k)=0\\
H_1&:\rho(k)\ne0\\
\left|\sqrt T\,\hat\rho(k)\right|&>1.96\Rightarrow \text{reject at approximately 5 percent}
\end{aligned}
\]

Variables: \(1.96\) is the two-sided 5 percent normal cutoff.

\subsubsection{Demeaned returns for LM predictability
tests}\label{demeaned-returns-for-lm-predictability-tests}

\[
e_t=r_t-\bar r
\]

Variables: \(\bar r\) is the sample mean return.

\subsubsection{LM return predictability
regression}\label{lm-return-predictability-regression}

\[
e_t=\gamma_0+\gamma_1e_{t-1}+\gamma_2e_{t-2}+\cdots+v_t
\]

Variables: lagged demeaned returns test predictability.

\subsubsection{LM return predictability hypotheses and
statistic}\label{lm-return-predictability-hypotheses-and-statistic}

\[
\begin{aligned}
H_0&:\gamma_1=\gamma_2=\cdots=\gamma_k=0\\
H_1&:\text{at least one }\gamma_j\text{ differs from zero}\\
AR(k)&=TR^2\sim\chi_k^2
\end{aligned}
\]

Variables: \(\gamma_j\) are lag coefficients; \(\chi_k^2\) is chi-square
with \(k\) degrees of freedom.

\subsubsection{Variance-ratio one-period log
return}\label{variance-ratio-one-period-log-return}

\[
r_t=\log(P_t)-\log(P_{t-1})
\]

Variables: \(r_t\) is one-period log return.

\subsubsection{Variance-ratio n-period log
return}\label{variance-ratio-n-period-log-return}

\[
r_t(n)=r_t+r_{t-1}+\cdots+r_{t-(n-1)}
\]

Variables: \(r_t(n)\) is the \(n\)-period log return.

\subsubsection{Variance scaling under
independence}\label{variance-scaling-under-independence}

\[
\operatorname{Var}[\text{n-period return}]=n\operatorname{Var}[\text{1-period return}]
\]

Variables: this holds approximately when returns are independent over
time.

\subsubsection{Variance ratio}\label{variance-ratio}

\[
VR_n=\frac{\text{variance of n-period returns}}{n\times\text{variance of 1-period returns}}
\]

Variables: \(VR_n=1\) suggests no autocorrelation; \(VR_n>1\) suggests
positive autocorrelation; \(VR_n<1\) suggests negative autocorrelation.

\end{document}
